% Part: first-order-logic
% Chapter: proof-systems
% Section: axiomatic-deduction

\documentclass[../../../include/open-logic-section]{subfiles}

\begin{document}

\iftag{FOL}
      {\olfileid{fol}{prf}{axd}}
      {\olfileid{pl}{prf}{axd}}

\olsection{স্বতঃসিদ্ধমূলক \usetoken{P}{derivation}}

স্বতঃসিদ্ধমূলক !!{derivation}s হলো যুক্তিবিদ্যার প্রাচীনতম ও সরলতম
!!{derivation} পদ্ধতি। এই পদ্ধতির !!{derivation}s নিছক
!!{sentence}s-এর অনুক্রম। !!{sentence}s-এর কোনো অনুক্রমকে সঠিক
!!{derivation} বলা হয়, যদি তার প্রতিটি !!{sentence}~$!A$ নিচের
কোনো একটি শর্ত পূরণ করে:
\begin{enumerate}
\item $!A$ একটি স্বতঃসিদ্ধ, অথবা
\item $!A$,~!!{sentence}s-এর প্রদত্ত সেট~$\Gamma$-র একটি !!a{element}, অথবা
\item কোনো অনুমান-বিধি দিয়ে $!A$ সমর্থিত।
\end{enumerate}
স্বতঃসিদ্ধ হতে হলে $!A$-কে কয়েকটি নির্দিষ্ট !!{sentence} ছকের কোনো
একটির রূপের হতে হবে। প্রথম-ক্রমের যুক্তিবিদ্যার জন্য সন্তোষজনক (বিশুদ্ধ ও
পূর্ণ) !!{derivation} পদ্ধতি দেয়, এমন বহু স্বতঃসিদ্ধ-ছকসমষ্টি আছে।
কয়েকটি সমষ্টি তাদের নিয়ন্ত্রিত সংযোজক অনুসারে সাজানো; যেমন,
\[
!A \lif (!B \lif !A) \qquad !B \lif (!B \lor !C) \qquad (!B \land !C) \lif !B
\]
ছকগুলি $\lif$, $\lor$ ও~$\land$ নিয়ন্ত্রণকারী প্রচলিত স্বতঃসিদ্ধ।
কিছু স্বতঃসিদ্ধ পদ্ধতির লক্ষ্য স্বতঃসিদ্ধের সংখ্যা যতটা সম্ভব কম রাখা।
কোন সংযোজকগুলি মৌলিক ধরা হচ্ছে তার উপর নির্ভর করে, এমনকি একটিমাত্র
স্বতঃসিদ্ধ নিয়ে গঠিত স্বতঃসিদ্ধ পদ্ধতিও পাওয়া সম্ভব।

অনুমান-বিধি হলো একটি শর্তাধীন বিবৃতি, যা কোনো !!a{derivation}-এ কোনো
!!a{sentence}-কে সমর্থিত ধাপ হিসেবে গণ্য করার যথেষ্ট শর্ত দেয়। মোডাস
পোনেন্স এই ধরনের অতি প্রচলিত একটি বিধি: এটি বলে, $!A$ ও $!A \lif
!B$ আগে থেকেই সমর্থিত হলে $!B$-ও সমর্থিত। অর্থাৎ কোনো
!!a{derivation}-এর~$!B$ !!{sentence}-যুক্ত একটি পংক্তি সমর্থিত,
যদি $!B$-এর আগে সেই !!{derivation}-এ $!A$ ও $!A \lif !B$—দুটিই
(কোনো !!{sentence}~$!A$-এর জন্য) থাকে।

স্বতঃসিদ্ধমূলক !!{derivation}s-ভিত্তিক $\Proves$ সম্পর্কটি নিচের মতো
সংজ্ঞায়িত: $\Gamma \Proves !A$ তখনই, যখন এমন একটি !!a{derivation}
আছে যার শেষ সূত্র !!{sentence}~$!A$ (আর উপরের~(2) দিয়ে সমর্থিত সেই
!!{derivation}-এর !!{sentence}s-এর সেট হিসেবে
$\Gamma$-কে নেওয়া হয়)। $!A$ তখনই একটি উপপাদ্য, যখন~$\Gamma$ শূন্য
এমন অবস্থায় $!A$-এর একটি !!a{derivation} থাকে; অর্থাৎ
!!{derivation}-এর প্রতিটি !!{sentence} হয় (1), নয়তো~(3) দিয়ে
সমর্থিত। যেমন, নিচের !!a{derivation}-টি দেখায় যে
$\Proves !A  \lif (!B \lif (!B \lor !A))$:
\begin{derivation}
  1. & $!B \lif (!B \lor !A)$ \\
  2. & $(!B \lif (!B \lor !A)) \lif (!A  \lif (!B \lif (!B \lor !A)))$\\
  3. & $!A  \lif (!B \lif (!B \lor !A))$
\end{derivation}
পংক্তি~1-এর !!{sentence}-টি $!A \lif (!A
\lor !B)$ স্বতঃসিদ্ধের রূপের ($!A$ ও $!B$-এর ভূমিকা অদলবদল করে)।
পংক্তি~2-এর বাক্যটি $!A \lif (!B \lif !A)$ স্বতঃসিদ্ধের রূপের। তাই
উভয় পংক্তিই সমর্থিত। পংক্তি~3 মোডাস পোনেন্স দিয়ে সমর্থিত:
সেটিকে সংক্ষেপে $!D$ লিখলে, পংক্তি~2-এর রূপ $!C \lif !D$, যেখানে
$!C$ হলো $!B \lif (!B \lor !A)$, অর্থাৎ পংক্তি~1।

$\Gamma \Proves \lfalse$ হলে কোনো সেট $\Gamma$ অসঙ্গত। একটি পূর্ণ
স্বতঃসিদ্ধ পদ্ধতি যেকোনো~$!A$-এর জন্য $\lfalse \lif !A$-ও প্রমাণ করবে;
তাই $\Gamma$ অসঙ্গত হলে যেকোনো~$!A$-এর জন্য $\Gamma \Proves !A$।

যুক্তিবিদ্যার স্বতঃসিদ্ধমূলক !!{derivation}s-এর পদ্ধতি প্রথম দেন
গটলব ফ্রেগে, তাঁর ১৮৭৯ সালের \emph{Begriffsschrift} গ্রন্থে; এই কারণেই
বইটিকে প্রায়ই আধুনিক যুক্তিবিদ্যার প্রথম রচনা বলা হয়। অ্যালফ্রেড নর্থ
হোয়াইটহেড ও বার্ট্রান্ড রাসেল তাঁদের \emph{Principia Mathematica}-য়,
এবং ডেভিড হিলবার্ট ও তাঁর ছাত্রেরা ১৯২০-এর দশকে এই পদ্ধতিগুলিকে
পরিণত রূপ দেন। তাই এগুলিকে প্রায়ই “ফ্রেগে পদ্ধতি” বা “হিলবার্ট
পদ্ধতি” বলা হয়। এগুলি খুব বহুমুখী, কারণ কোনো যুক্তিবিদ্যার জন্য
স্বতঃসিদ্ধমূলক পদ্ধতি খুঁজে পাওয়া প্রায়ই সহজ। !!{derivation}s-এর গঠন
অত্যন্ত সরল এবং অনুমান-বিধি মাত্র একটি বা দুটি বলে, তাদের \emph{সম্বন্ধে}
বিভিন্ন বিষয় প্রমাণ করাও তুলনামূলকভাবে সহজ। তবে কার্যক্ষেত্রে এগুলি
ব্যবহার করা খুব কঠিন; অর্থাৎ প্রমাণ খুঁজে পাওয়া ও লেখা কঠিন।

\end{document}
